% Irodov Part III §3.3 113--121; source Fig 3.18--3.23 inspected,
% user PDF pp.116--117 (printed pp.107--108).
\electrocapproblem{113}{107}
图3.18中四块面积相同为 $S$ 的平行导体板，
自上至下相邻间距都是 $d$。按图(a)、(b)
所示导线连接后，求 A、B 间电容。
\exsolution
记每个相邻间隙的平板电容
$C_0=\varepsilon_0S/d$。
(a) 图(a)把第1、3板通过左端导线短接，
第2板接 A，第4板接 B。
第1—2和第2—3间两个 $C_0$
并联成 $2C_0$，它再与第3—4间的
$C_0$ 串联，故
\[
\boxed{C_a=
\frac{(2C_0)C_0}{2C_0+C_0}
=\frac{2\varepsilon_0S}{3d}}.
\]
(b) 图(b)把第1、4板通过左端导线短接，
第2板接 A，第3板接 B。
第1—2间 $C_0$ 和第3—4间 $C_0$
以中间浮置导体节点形成串联支路，
等效 $C_0/2$；同时第2—3间
$C_0$ 直接跨 A、B 并联，故
\[
\boxed{C_b=C_0+\frac{C_0}{2}
=\frac{3\varepsilon_0S}{2d}}.
\]
\excheck 本题需要先按图识别短接的
\emph{同势金属板}，才能进行串并联化简；
不能把四块板按几何顺序简单串联。

\electrocapproblem{114}{107}
两电容器 $C_1=1.0\,\mathrm{\mu F}$，
最大耐压 $V_1=6.0\,\mathrm{kV}$；
$C_2=2.0\,\mathrm{\mu F}$，
最大耐压 $V_2=4.0\,\mathrm{kV}$。
串联后的最大安全总电压是多少？
\exsolution
串联时两电容器携带相同电荷 $Q$，
各自电压 $U_i=Q/C_i$。
安全条件为
$Q\le C_1V_1$ 与 $Q\le C_2V_2$，
所以
\[
Q_{\max}=\min(1.0\times6.0,\,
2.0\times4.0)\,\mathrm{mC}
=6.0\,\mathrm{mC}.
\]
总电压为
\[
V_{\max}=Q_{\max}
\left(\frac1{C_1}+\frac1{C_2}\right)
=V_1\left(1+\frac{C_1}{C_2}\right).
\]
因此
\[
\boxed{V_{\max}=9.0\,\mathrm{kV}}.
\]
\excheck 先到耐压极限的是 $C_1$；
此时 $C_2$ 电压为3 kV，低于4 kV。

\electrocapproblem{115}{107}
原书图3.19：一个电动势
$\mathcal E=110\,\mathrm V$ 的电源，
上支路两个相同电容 $C_1$ 串联，
各串联节点到下导线 B 另接电容 $C_2$；
末端节点为 A。
已知 $C_2/C_1=\eta=2.0$。
求 $\varphi_A-\varphi_B$。
\exsolution
取下端 B 电势为零，
电源上端电势 $\mathcal E$。
把中间节点电势记为 $x$，
末端 A 电势为 $y$。
未由外部注入自由电荷的浮置节点需净电荷为零：
\[
C_1(x-\mathcal E)+C_1(x-y)+C_2x=0,
\]
\[
C_1(y-x)+C_2y=0.
\]
用 $\eta=C_2/C_1$ 化简：
$(2+\eta)x-y=\mathcal E$、
$x=(1+\eta)y$。
相消得到
\[
\boxed{y=\varphi_A-\varphi_B=
\frac{\mathcal E}{1+3\eta+\eta^2}}.
\]
代入 $\eta=2$：
\[
\boxed{\varphi_A-\varphi_B=
\frac{110}{11}\,\mathrm V=10\,\mathrm V}.
\]
\excheck 若误把两个 $C_2$ 直接并联，
会破坏中间节点电荷守恒，
无法得到该分母。

\electrocapproblem{116}{107}
图3.20：一个从左向右无限重复的电容梯形网络，
每节沿上支路有电容 $C$，
节点到公共下支路又有电容 $C$。
求左端与下公共导线之间的总等效电容。
\exsolution
设整个无穷网络的等效电容为 $X$。
第一只上支路电容 $C$ 后，
从新节点向下有一个电容 $C$，
并联上与原网络形状完全相同的后续部分，
其电容亦为 $X$，故节点到地等效为 $C+X$。
因此
\[
\frac1X=\frac1C+\frac1{C+X}.
\]
令 $x=X/C$，整理：
$x=(1+x)/(2+x)$，
即
$x^2+x-1=0$。
正的物理解为
\[
\boxed{X=\frac{\sqrt5-1}{2}C\simeq0.618C}.
\]
\excheck 必须取正根；
另一个根 $-(1+\sqrt5)/2$ 不能为被动电容的等效值。

\electrocapproblem{117}{107--108}
图3.21的 A、B 间依次串接
电容器 $C_1=1.0\,\mathrm{\mu F}$、
电源 $\mathcal E=10\,\mathrm V$、
电容器 $C_2=2.0\,\mathrm{\mu F}$。
已知 $\varphi_A-\varphi_B=5.0\,\mathrm V$。
求两只电容器各自的电压大小。
\exsolution
两只电容器串联而中间电源可施加电势跃变，
两只电容器的电荷模相同，设为 $Q$。
按原书图3.21的电源极性，
两电容器两端电压和
$U_1+U_2=(\varphi_A-\varphi_B)+\mathcal E$。
因此
\[
Q\left(\frac1{C_1}+\frac1{C_2}\right)
=5.0+10.0=15.0\,\mathrm V,
\]
\[
Q=\frac{(15.0)C_1C_2}{C_1+C_2}
=10\,\mathrm{\mu C}.
\]
所以
\[
\boxed{U_1=\frac Q{C_1}=10\,\mathrm V},
\qquad
\boxed{U_2=\frac Q{C_2}=5.0\,\mathrm V}.
\]
\excheck 带电状态与电源极性相关；
电压\emph{模}为正，代数电压沿回路应保留方向。

\electrocapproblem{118}{108}
图3.22：一个闭合环路中两只电容器
$C_1,C_2$ 与两个电动势
$\mathcal E_1,\mathcal E_2$ 按图中方向连接。
求每只电容器左极板相对于右极板的电势差。
\exsolution
设电容器极板代数电荷模为 $Q$，
在无旁路充电的串联闭环中各自绝对值相同。
沿图3.22所取的电动势方向，
两个电源的代数贡献为
$\mathcal E_2-\mathcal E_1$，
所以
\[
Q\left(\frac1{C_1}+\frac1{C_2}\right)
=\mathcal E_2-\mathcal E_1.
\]
对两个电容器左至右的电压，
依图中板极性设定：
\[
\boxed{V_1=\frac{\mathcal E_2-\mathcal E_1}
{1+C_1/C_2}},\qquad
\boxed{V_2=\frac{\mathcal E_1-\mathcal E_2}
{1+C_2/C_1}}.
\]
\excheck 两只电容器的左—右电压并非
同时同号；若交换源的正负端，
上述两个式子的代数符号须相应改变。

\electrocapproblem{119}{108}
仍在图3.22的电源与电容串联闭环中，
求每只电容器上的电荷\emph{模}。
\exsolution
由图中两个电源的代数电动势差
$\Delta\mathcal E=\mathcal E_2-\mathcal E_1$，
稳态的电容回路平衡为
$Q/C_1+Q/C_2=|\Delta\mathcal E|$。
于是
\[
\boxed{Q_1=Q_2=
\frac{C_1C_2}{C_1+C_2}
|\mathcal E_2-\mathcal E_1|}.
\]
\excheck 两只电容器带电量大小相等，
且 $\mathcal E_1=\mathcal E_2$ 时为零；
这是在未预充电、无其它净节点电荷条件下的结果。

\electrocapproblem{120}{108}
图3.23的四臂电容桥，
各臂电容为 $C_1,C_2,C_3,C_4$，
上下接电动势 $\mathcal E$，
取两个中间节点为 A、B。
求 $\varphi_A-\varphi_B$，以及零电势差条件。
\exsolution
电容桥两条并联支路分别为
$C_1$ 与 $C_3$ 串联、
$C_2$ 与 $C_4$ 串联；
选图中上下源极性，
在支路的电压分配公式中
\[
\varphi_A=\frac{C_2}{C_1+C_2}\mathcal E
\quad\text{与}\quad
\varphi_B=\frac{C_4}{C_3+C_4}\mathcal E
\]
在原书节点标号约定下，
对应两段电势差整理为
\[
\boxed{\varphi_A-\varphi_B=
\mathcal E\frac{C_2C_3-C_1C_4}
{(C_1+C_2)(C_3+C_4)}}.
\]
桥平衡条件是分子为零：
\[
\boxed{C_1/C_2=C_3/C_4}.
\]
\excheck 图中四个电容的编号不是
顺序围绕一个矩形而不变；
应按每一串联支路中
浮动节点的电荷守恒确定电压。

\electrocapproblem{121}{108}
一只已充至 $V=110\,\mathrm V$ 的
电容器 $C_1=1.0\,\mathrm{\mu F}$，
并联接到另外两只原先不带电、串联的
电容器 $C_2=2.0\,\mathrm{\mu F}$、
$C_3=3.0\,\mathrm{\mu F}$ 的两端。
求通过连接导线的电荷量。
\exsolution
设两只原先未带电的串联电容器最终带电荷模
$q$，其所分配电压为
$V_f=q(1/C_2+1/C_3)$。
$C_1$ 最终带电量为 $C_1V_f$，
而整个外部节点总电荷守恒：
$C_1V=C_1V_f+q$。
联立可得
\[
C_1V=q\left[1+C_1
\left(\frac1{C_2}+\frac1{C_3}\right)\right],
\]
即
\[
\boxed{q=\frac{V}
{1/C_1+1/C_2+1/C_3}}.
\]
代入数值：
\[
\boxed{q=\frac{110}{1+1/2+1/3}\,\mathrm{\mu C}
=60\,\mathrm{\mu C}=0.060\,\mathrm{mC}}.
\]
\excheck 电容值以 $\mu$F 计时，
式中 $V/(\mu\mathrm F^{-1})$ 为 $\mu$C；
通过任一串联连接导线的累计电荷量相同。
